\section{Conclusion}
\label{sec:conclusion}

We built and verified an open optimisation model of a hydropower-dependent electricity system. It couples
hourly dispatch, a physical reservoir water balance, batteries, imports and investment under uncertainty,
and it reports every result together with its data provenance and its limitations. Under the stated (partly
synthetic) assumptions, the experiments support the following defensible findings.
\begin{enumerate}
\item Drought is primarily a financial and import-dependence risk. Reliability is governed by the
availability of interconnection: shortages require import availability of about 50\% or less together with
severe drought or demand growth.
\item Solar PV is a no-regret investment in every scenario as long as its surplus can be exported. Its
marginal value is about 1.7$\times$ its annualised cost, and the technology choice is robust to every
sensitivity tested.
\item Batteries do not recover their cost through arbitrage, and they increase gross imports. They do pay
off through avoided unserved energy in joint drought and import-stress events. Discharge power is more
valuable than additional energy capacity.
\item Stochastic, CVaR-based planning changes the decision only when the dominant resource saturates. With
exports restricted, it buys a wind hedge that lowers tail cost by 74.5~MEUR/yr for 15.1~MEUR/yr of expected
cost, and this is confirmed by full chronology.
\item Methodologically, temporal resolution and the boundary treatment of representative periods for
seasonal storage each changed conclusions. Water-value boundaries and chronological validation are
inexpensive safeguards.
\end{enumerate}

\paragraph{Next steps.} (i)~Replace the synthetic demand and prices with observed ENTSO-E load and regional
market data. (ii)~Calibrate the hydrology against gauge or plant records and represent the three Drin stages
separately, with variable head. (iii)~Add network constraints for the four interconnections and correlated
regional scenarios, including the low-availability regime in which storage becomes valuable. (iv)~Add
reserve requirements and multi-stage hydro scheduling, for example SDDP-style water-value cuts. (v)~Estimate
scenario probabilities from climate-model ensembles rather than by assumption.
